\documentclass{article}

\usepackage[dblblindworkshop]{neurips_2026}
\workshoptitle{Sim2Science}

\usepackage[utf8]{inputenc}
% U+02BC (modifier letter apostrophe) appears once in audit/BIBLIOGRAPHY.bib, in
% Crossref's own fetched record for Brynjarsdottir & O'Hagan (2014) -- not in
% inputenc's default utf8 table. Declared here rather than edited in the .bib,
% which stays exactly as fetched (audit/BIBLIOGRAPHY.bib's own stated policy).
\DeclareUnicodeCharacter{02BC}{\textquoteright}
\usepackage[T1]{fontenc}
\usepackage{hyperref}
\usepackage{url}
\usepackage{booktabs}
\usepackage{amsfonts}
\usepackage{amsmath}
\usepackage{nicefrac}
\usepackage{microtype}
\usepackage{xcolor}
\usepackage{graphicx}

% Title as operator-confirmed. FLAGGED (not silently kept, not silently changed) in
% audit/S8_REPORT.md P-2: the C1 amendment removes any selective-inference METHOD
% claim, and whether "Selective Inference" as the lead term still fits the
% diagnostics-and-limits reframing is an open question for the operator, not a
% session. See the report for the alternative proposed there.
\title{Selective Inference for Component-Level Simulator Misspecification}

\author{%
  Anonymous Author(s) \\
  Anonymous Institution \\
  Anonymous Address \\
  \texttt{anonymous@example.com} \\
}

\begin{document}

\maketitle

\begin{abstract}
TODO -- drafted last, Phase 5, after every section below exists to preview accurately.
\end{abstract}

\section{Introduction}
\label{sec:intro}
TODO -- drafted last among body sections (Phase 3), once the actual content exists to
preview. Must motivate the attribution problem and preview the four findings under the
diagnostics-and-limits framing of the C1 amendment (\texttt{audit/FINAL\_CLAIMS.md}), not
as a methods paper.

\section{Background and related work}
\label{sec:background}

\paragraph{Detecting misspecification versus attributing it.}
Recent work in simulation-based inference (SBI) has moved from asking whether a simulator is
misspecified to asking where. \citet{montel2025testsmodelmisspecificationsimulationbased}
give a family of tests, from local distortions to global model checks, for detecting
misspecification of an SBI posterior; their global test corrects for multiplicity by
computing a trials-corrected $p$-value on the \emph{minimum} of per-summary statistics over
the candidate distortions, accounting for the correlations between them. This controls the
global null hypothesis that the simulator is correctly specified \emph{somewhere}; it does
not control the identity of the arg\,min, and their own presentation is framed throughout as
detection rather than as attribution to a specific simulator component. We take this as the
baseline our diagnostic sits downstream of: before any attribution procedure -- selective,
competitive, or otherwise -- is run at all, it is worth asking whether the summaries in hand
can distinguish the components being attributed to.

\paragraph{The identifiability precondition.}
That question has an exact answer in a closely related setting.
\citet{Kahl_2019} show that per-component input reconstruction in a dynamic system is
identifiable if and only if a structural sensitivity map has full column rank, resting on the
invertibility condition of \citet{Sain_1969}; \citet{CatchpoleMorgan_1997} give the general
parameter-redundancy detection framework a rank condition of this kind instantiates, and
\citet{Brynjarsd_ttir_2014} frame the broader identifiability problem for model discrepancy in
computer-model calibration. We treat all four as an established precondition, not as a result
of this paper: the diagnostic we use decides whether that precondition holds for a given
simulator and summary set, it does not establish that the precondition is the right one to
check.

\paragraph{Diagnostic tooling, applied rather than proposed.}
Screening a sensitivity matrix's rank and condition number to decide whether a set of
parameters is jointly recoverable is not new; \citet{Cintr_n_Arias_2009} apply exactly this
screen to a compartmental epidemic model, the same class of simulator we use in
Section~\ref{sec:experiments}. Because our sensitivity matrix is estimated by simulation
rather than known symbolically, we also need to know how far a finite-difference estimate can
be trusted under simulation noise; \citet{Mor__2011} and \citet{Mor__2012} give the relevant
noise-estimation and near-optimal-step-size results, of which our own step-size sweep
(Section~\ref{sec:method}) is a coarser, pre-registered substitute. Where the resulting
spectrum sits -- whether it has a clean gap or decays smoothly across many decades, in which
case any rank tolerance is partly a statement about where the analyst put it -- is the
objection \citet{Gutenkunst_2007} raise for sensitivity spectra in systems-biology models
generally; Section~\ref{sec:experiments} measures rather than assumes an answer for our own
simulator.

\paragraph{Exact inference under a composite, simulation-only null.}
Section~\ref{sec:negative-result} evaluates a natural construction for turning the diagnostic's
output into a calibrated test: exact finite-sample inference under a null distribution that
depends on unknown nuisance parameters, by maximizing a simulated $p$-value over the nuisance
set, is \citet{Dufour_2006}'s maximized Monte Carlo test; calibrating a selection event that
has no analytic description, by rejection-sampling the simulator's own null until a replicate
lands in the observed selection cell, is available for arbitrary conditioning events following
\citet{freidling2026selectiverandomizationinferenceadaptive}. Composing the two lets a
simulator supply everything both constructions need without any further modelling assumption;
we cite both as established tools and evaluate the composition on our own simulator, rather
than propose either as new.

\section{Method}
\label{sec:method}

\paragraph{Scope assumption, stated first.}
Throughout this paper we assume each simulator component is misspecified by at most one
mechanism at a time: every component carries a single one-parameter distortion family, never
two at once. This assumption is load-bearing rather than a footnote -- Section~\ref{sec:experiments}
shows directly what happens once it is dropped -- so we state it here, before any of the
diagnostic's machinery, rather than after the results that depend on it.

\paragraph{The diagnostic.}
Every ingredient below is prior art, applied rather than proposed (Section~\ref{sec:background}).
Declare the simulator as $K$ components, each carrying a one-parameter distortion family
$\delta_k(\cdot\,;\eta_k)$ that is bit-identical to the base simulator at $\eta_k = 0$. Estimate
the summary Jacobian $J = \partial s/\partial\eta$ at $\eta=0$ by central differences under
common random numbers, swept over six decades of step size $h$; without common random numbers
the finite difference is dominated by simulation noise of order $\sigma/h$ and no stable range
of $h$ exists, which we confirm with a negative control. Summaries are divided by their
prior-predictive standard deviation and each $\eta_k$ by a fixed relative perturbation scale --
both normalizations, and every tolerance below, fixed before any singular value in this paper
existed. The set of components is \emph{separable} under a summary set if $J$ has full column
rank at tolerance $\tau$ and condition number $\kappa \le \kappa_{\max}$; we use $\tau=10^{-2}$
and $\kappa_{\max}=1/\tau=100$, chosen from a compute budget rather than a convention -- resolving
components to condition number $\kappa$ costs on the order of $\kappa^2$ simulator replicates,
and $\kappa=100$ is where that cost exhausts the budget available for this study. A singular
value that varies by more than a factor of two across the step-size plateau is \emph{unresolved}
and is counted toward the rank in neither direction, so the diagnostic can report ``not
determined at this precision'' as a third outcome distinct from separable and inseparable.

\paragraph{Where separability fails.}
When the rank condition fails, we report the \emph{equivalence class} of components named by
each near-null right singular vector $v$ ($\sigma \le \tau\sigma_1$), rather than naming a single
culprit: component $k$ joins the class if $|v_k|\ge 0.3$, a threshold set below both the
two-component confound value ($|v_k|\approx 0.707$) and the three-component one ($\approx 0.577$)
by a comfortable margin. Pairwise coherence is reported alongside for interpretation but is not
the decision rule: the coherence consistent with $\kappa_{\max}=100$ is $\mu\approx0.9998$, too
close to one to be estimable, so we use it only to point at which pair of columns is responsible
for a flagged near-null direction. No arrow leaves the verdict -- the diagnostic decides whether
attribution is well posed for a given simulator and summary set, and nothing further.

\section{Experiments}
\label{sec:experiments}

\paragraph{Simulator and summary sets.}
We apply the diagnostic to a $K=3$ compartmental epidemic simulator (Figure~\ref{fig:simulator}):
transmission, progression, and observation, each carrying one declared distortion family
that is the base simulator exactly at $\eta_k=0$, and two declared family sets (base and
adversarial, the latter designed so all three columns perturb the same feature -- the epidemic
curve's exponential growth rate). We evaluate three summary sets: $S_A$, epidemic-curve
features a domain modeller would choose (peak height, peak time, final size, growth rate);
$S_B$, roughly ten time-binned incidence counts an SBI practitioner would choose; and $S_C$, an
impoverished positive control with only two summaries, which \emph{must} fail by construction
($d=2<K=3$) and does -- rank 2, $\kappa=\infty$ -- confirming the diagnostic reports genuine
rank deficiency rather than an artifact of the harness. A random-attributor floor of $1/K$
(analytic $0.3333$, measured $0.3299$ over $10{,}000$ draws) is reported against every accuracy
figure elsewhere in this line of work.

\paragraph{Where attribution is identifiable: $S_B$ under eight family assignments.}
Combining the two declared families component-wise gives eight distinct three-column
distortion models (Figure~\ref{fig:assignments}), only two of which -- all-base and
all-adversarial -- were the ones originally declared; the other six recombine the same
columns and cost no additional simulation. Under our scope assumption, $S_B$ separates the
three components under \emph{all eight}: condition number from $6.628$ to $65.64$, every
singular value resolved, no coherence pair flagged. $S_A$ is a control rather than a
result: it splits exactly on the transmission family, separable under every assignment that
keeps the base transmission distortion and inseparable under every one that swaps in the
adversarial one, failing at $\kappa=100.9$ -- $0.9\%$ past the ceiling -- on the assignment
that changes only that one family. An instrument that always reported ``separable'' could not
produce that contrast. The margin is not uniform: $S_B$'s verdict survives $\tau$ moving by a
factor of $9.881$ under the base assignment but only $1.523$ under the tightest one
(Figure~\ref{fig:threshold}), so doubling $\tau$ -- the coarse robustness check this project's
own protocol proposes -- flips it.

\paragraph{Where it is not: two distortion parameters per component.}
Placing all six declared distortion directions side by side -- two per component -- gives a
single six-column object (there is only one; the two family sets do not define two of them).
For $S_B$ it is \textbf{inseparable} at $\kappa=628.9$, rank $4$ of $6$
(Figure~\ref{fig:spectrum}). This is not an unresolved estimator (worst singular-value
variation across the step-size plateau: $1.023$, against an admissible $2$); not a threshold
artifact (inseparable across $\tau$ from $0.005$ to $1.0$, and there is no tolerance at which
the object is both separable and affordable); and not structural ($d=10\ge6$, so no singular
value is zero by construction). Both near-null directions are cross-mechanism, naming
progression together with observation, with transmission carrying under $5\%$ of the energy in
either (Figure~\ref{fig:confound}): a drifting removal hazard is nearly indistinguishable from
a constant hazard change combined with a drifting reporting rate. The confound needs the
progression component to carry two distortion parameters at once, which is exactly why it
cannot arise under our scope assumption's one-parameter-per-component restriction, and exactly
why that restriction -- not a bound on $\kappa$ -- is what the assumption has to be. At $K=6$
the spectrum also matches the gapless, no-clean-break picture the sloppy-models literature
describes \citep{Gutenkunst_2007} (spread $2.80$ decades, gap prominence $1.40$, $\tau\sigma_1$
falling inside the spectrum); at $K=3$ it does not. The reassurance a $K=3$ rank threshold
offers is a fact about $K=3$, not about the simulator.

The full eight-assignment table for both summary sets, and the complete threshold-sensitivity
sweep, are reproduced in Appendix~\ref{sec:appendix}.

\section{The MMC negative result}
\label{sec:negative-result}
TODO -- Phase 3. The boundary sweep (Fig.~\ref{fig:nontermination}), the nuisance-to-noise
ratio predicted before measurement and confirmed by it. Stated as a genuine finding, not
an apology for an unbuilt experiment.

\section{Limitations}
\label{sec:limitations}
TODO -- Phase 3. D-14's scope assumption and what it excludes; R2's NARROW-CONDITIONAL
status where load-bearing; single-simulator generalization; the three hand-annotated
figures' annotations outside the automated provenance check.

\section*{Reproducibility checklist}
\input{checklist.tex}

\bibliographystyle{plainnat}
\bibliography{../audit/BIBLIOGRAPHY}

\appendix
\section{Appendix}
\label{sec:appendix}
TODO -- Phase 2--4, as needed. Full 8-assignment table, threshold-sensitivity sweep
detail, \texttt{audit/FINAL\_CLAIMS.md}'s generated-numbers appendix reproduced verbatim
(S11).

\end{document}
